\documentclass{beamer}

\usetheme{Madrid}

\usepackage{amsmath,amsfonts,amssymb,graphicx}

\title{Recurrent Neural Networks}
\subtitle{Mathematics, Algorithms, and Applications}
\author{Morten Hjorth-Jensen}
\date{FYS5429/9429 March 5}

\begin{document}

\frame{\titlepage}

%------------------------------------------------
\section{Motivation}

\begin{frame}{Why Sequential Models?}

Many physical and computational problems involve sequences:

\begin{itemize}
\item time series analysis
\item natural language processing
\item dynamical systems
\item signal processing
\item control problems
\end{itemize}

Examples of applications in the physical sciences:

\begin{itemize}
\item molecular dynamics trajectories
\item quantum time evolution
\item climate modelling
\item detector signals
\end{itemize}

Standard feed-forward networks assume independent samples.
\end{frame}

%------------------------------------------------
\begin{frame}{Sequential Data Structure}

Instead of a single input vector:

\[
x \in \mathbb{R}^d
\]

we have a sequence

\[
x_1,x_2,\dots,x_T
\]

Each element:

\[
x_t \in \mathbb{R}^d
\]

represents the system at time step \(t\).
\end{frame}

%------------------------------------------------
\section{Recurrent Neural Networks}

\begin{frame}{Core Idea}

Recurrent neural networks introduce a hidden state

\[
h_t
\]

that summarizes the past.

\[
h_t = f(x_t,h_{t-1})
\]

The hidden state acts as a memory.
\end{frame}

%------------------------------------------------
\begin{frame}{RNN Update Equations}

Standard formulation:

\[
h_t = \sigma(W_{xh}x_t + W_{hh}h_{t-1} + b_h)
\]

Output:

\[
y_t = \phi(W_{hy}h_t + b_y)
\]

Parameters:

\[
\theta = \{W_{xh},W_{hh},W_{hy},b_h,b_y\}
\]
\end{frame}

%------------------------------------------------
\section{Forward Propagation}

\begin{frame}{Forward Propagation}

For each time step:

\[
h_t = \sigma(W_{xh}x_t + W_{hh}h_{t-1} + b_h)
\]

\[
y_t = \phi(W_{hy}h_t + b_y)
\]

Initial state:

\[
h_0 = 0
\]

or learned parameters.
\end{frame}

%------------------------------------------------
\begin{frame}{Unrolling the Network}

RNNs can be viewed as deep networks over time.

\[
x_1 \rightarrow h_1 \rightarrow y_1
\]

\[
x_2 \rightarrow h_2 \rightarrow y_2
\]

\[
x_3 \rightarrow h_3 \rightarrow y_3
\]

Weights are shared across time.
\end{frame}

%------------------------------------------------
\section{Loss Functions}

\begin{frame}{Sequence Loss}

Define total loss:

\[
L = \sum_{t=1}^{T} \ell(y_t,\hat y_t)
\]

Examples:

Mean squared error:

\[
\ell = ||y_t-\hat y_t||^2
\]

Cross entropy:

\[
\ell = - \sum_i \hat y_i \log y_i
\]
\end{frame}

%------------------------------------------------
\section{Backpropagation Through Time}

\begin{frame}{Training the Network}

Training requires computing

\[
\frac{\partial L}{\partial \theta}
\]

Since hidden states depend on earlier states, gradients must propagate through time.

This algorithm is called

\textbf{Backpropagation Through Time (BPTT)}.
\end{frame}

%------------------------------------------------
\begin{frame}{Hidden State Gradient}

Define

\[
\delta_t = \frac{\partial L}{\partial h_t}
\]

Using the chain rule:

\[
\delta_t =
\frac{\partial \ell_t}{\partial h_t}
+
W_{hh}^T\delta_{t+1}
\]
\end{frame}

%------------------------------------------------
\begin{frame}{Weight Gradients}

\[
\frac{\partial L}{\partial W_{hh}}
=
\sum_{t=1}^T \delta_t h_{t-1}^T
\]

\[
\frac{\partial L}{\partial W_{xh}}
=
\sum_{t=1}^T \delta_t x_t^T
\]

\[
\frac{\partial L}{\partial W_{hy}}
=
\sum_{t=1}^T
\frac{\partial \ell_t}{\partial y_t}
h_t^T
\]
\end{frame}

%------------------------------------------------
\section{Gradient Instabilities}

\begin{frame}{Vanishing Gradients}

Backpropagation involves repeated multiplication:

\[
(W_{hh})^k
\]

If eigenvalues satisfy

\[
|\lambda| < 1
\]

gradients decay exponentially.
\end{frame}

%------------------------------------------------
\begin{frame}{Exploding Gradients}

If

\[
|\lambda| > 1
\]

gradients grow exponentially.

Solutions:

\begin{itemize}
\item gradient clipping
\item improved architectures
\item careful initialization
\end{itemize}
\end{frame}


%=========================================================
\section{BPTT: Full Derivation with Matrix Calculus}
%=========================================================

\begin{frame}{Setup and Notation (Vectorized)}
  We consider a (single-sequence) vanilla RNN with hidden size \(m\), input size \(d\), output size \(k\).

  Forward equations:
  \[
  a_t = W_{xh} x_t + W_{hh} h_{t-1} + b_h,\qquad h_t = \sigma(a_t)
  \]
  \[
  o_t = W_{hy} h_t + b_y,\qquad y_t = \phi(o_t)
  \]

  Dimensions:
  \[
  W_{xh}\in\mathbb{R}^{m\times d},\quad W_{hh}\in\mathbb{R}^{m\times m},\quad b_h\in\mathbb{R}^m
  \]
  \[
  W_{hy}\in\mathbb{R}^{k\times m},\quad b_y\in\mathbb{R}^k
  \]

  Total loss (sequence-to-sequence):
  \[
  L=\sum_{t=1}^T \ell_t(y_t,\hat y_t).
  \]
\end{frame}

%---------------------------------------------------------
\begin{frame}{Matrix Calculus Rules Used}
  We use the following standard identities:

  \begin{itemize}
  \item For \(z = Ax\), \( \frac{\partial L}{\partial x} = A^T \frac{\partial L}{\partial z}\).
  \item For \(z = Ax + b\), \( \frac{\partial L}{\partial A} = \frac{\partial L}{\partial z}\, x^T\), \(\frac{\partial L}{\partial b} = \frac{\partial L}{\partial z}\).
  \item For elementwise nonlinearity \(h=\sigma(a)\):
    \[
    \frac{\partial L}{\partial a} = \frac{\partial L}{\partial h}\odot \sigma'(a),
    \]
    where \(\odot\) is Hadamard product.
  \end{itemize}

  We will compute gradients by backpropagating through the unrolled computational graph.
\end{frame}

%---------------------------------------------------------
\begin{frame}{Local Output Gradients}
  Define output-layer ``error signals'':
  \[
  \delta^o_t \equiv \frac{\partial L}{\partial o_t}\in\mathbb{R}^k.
  \]

  By chain rule:
  \[
  \delta^o_t = \left(\frac{\partial \ell_t}{\partial y_t}\right)\odot \phi'(o_t).
  \]

  Then (matrix calculus):
  \[
  \frac{\partial L}{\partial W_{hy}} = \sum_{t=1}^T \delta^o_t\, h_t^T,
  \qquad
  \frac{\partial L}{\partial b_y} = \sum_{t=1}^T \delta^o_t.
  \]

  Hidden contribution from output at time \(t\):
  \[
  \left.\frac{\partial L}{\partial h_t}\right|_{\text{out}}
  =
  W_{hy}^T \delta^o_t.
  \]
\end{frame}

%---------------------------------------------------------
\begin{frame}{Key Quantity: Hidden-State Sensitivity}
  We want the total sensitivity
  \[
  \delta^h_t \equiv \frac{\partial L}{\partial h_t}\in\mathbb{R}^m.
  \]

  Because \(h_t\) influences:
  \begin{itemize}
  \item the loss at time \(t\) via \(o_t\),
  \item all future hidden states \(h_{t+1},h_{t+2},\dots\) through recurrence.
  \end{itemize}

  Thus:
  \[
  \delta^h_t
  =
  \left.\frac{\partial L}{\partial h_t}\right|_{\text{out}}
  +
  \left.\frac{\partial L}{\partial h_t}\right|_{\text{future}}.
  \]

  We now compute the future term via the recurrence Jacobian.
\end{frame}

%---------------------------------------------------------
\begin{frame}{From \(\delta^h_t\) to \(\delta^a_t\)}
  Define the pre-activation error:
  \[
  \delta^a_t \equiv \frac{\partial L}{\partial a_t}\in\mathbb{R}^m.
  \]

  Since \(h_t=\sigma(a_t)\) elementwise:
  \[
  \delta^a_t = \delta^h_t \odot \sigma'(a_t).
  \]

  This \(\delta^a_t\) is the central object used to compute parameter gradients in the recurrent core.
\end{frame}

%---------------------------------------------------------
\begin{frame}{Recurrence Jacobian and Backward Recursion}
  We have
  \[
  a_{t+1} = W_{xh}x_{t+1}+W_{hh}h_t+b_h.
  \]

  Hence the Jacobian of \(a_{t+1}\) w.r.t. \(h_t\) is:
  \[
  \frac{\partial a_{t+1}}{\partial h_t} = W_{hh}.
  \]

  Using matrix calculus:
  \[
  \left.\frac{\partial L}{\partial h_t}\right|_{\text{future}}
  =
  \left(\frac{\partial a_{t+1}}{\partial h_t}\right)^T
  \frac{\partial L}{\partial a_{t+1}}
  =
  W_{hh}^T \delta^a_{t+1}.
  \]

  Therefore the BPTT recursion is:
  \[
  \boxed{\delta^h_t = W_{hy}^T\delta^o_t + W_{hh}^T\delta^a_{t+1}},
  \qquad
  \boxed{\delta^a_t = \delta^h_t \odot \sigma'(a_t)}
  \]
  with terminal condition \(\delta^a_{T+1}=0\).
\end{frame}

%---------------------------------------------------------
\begin{frame}{Core Parameter Gradients via Matrix Calculus}
  Recall:
  \[
  a_t = W_{xh}x_t + W_{hh}h_{t-1} + b_h.
  \]

  Using the identity \(\partial L/\partial A = (\partial L/\partial z)\,x^T\) for \(z=Ax\):

  \[
  \boxed{
    \frac{\partial L}{\partial W_{xh}} = \sum_{t=1}^T \delta^a_t\, x_t^T}
  \]

  \[
  \boxed{
    \frac{\partial L}{\partial W_{hh}} = \sum_{t=1}^T \delta^a_t\, h_{t-1}^T}
  \]

  \[
  \boxed{
    \frac{\partial L}{\partial b_h} = \sum_{t=1}^T \delta^a_t}
  \]

  These are the standard BPTT gradients for a vanilla RNN.
\end{frame}

%---------------------------------------------------------
\begin{frame}{Explicit ``Unrolled'' Form (Useful for Analysis)}
  Iterating the recursion:
  \[
  \delta^h_t = W_{hy}^T\delta^o_t + W_{hh}^T(\delta^h_{t+1}\odot\sigma'(a_{t+1})).
  \]

  Expanding gives:
  \[
  \delta^h_t
  =
  W_{hy}^T\delta^o_t
  +
  W_{hh}^T D_{t+1} W_{hy}^T\delta^o_{t+1}
  +
  W_{hh}^T D_{t+1} W_{hh}^T D_{t+2} W_{hy}^T\delta^o_{t+2}
  +\cdots
  \]
  where \(D_t=\mathrm{diag}(\sigma'(a_t))\).

  This makes the role of repeated Jacobian products explicit.
\end{frame}

%---------------------------------------------------------
\begin{frame}{Vanishing/Exploding Gradients from Jacobian Products}
  The ``future'' influence involves products like:
  \[
  (W_{hh}^T D_{t+1})(W_{hh}^T D_{t+2})\cdots (W_{hh}^T D_{t+k}).
  \]

  If the operator norm satisfies (roughly)
  \[
  \|W_{hh}\|\cdot \max_t\|D_t\| < 1,
  \]
  then gradients shrink exponentially (vanishing gradients).

  If
  \[
  \|W_{hh}\|\cdot \max_t\|D_t\| > 1,
  \]
  then gradients can grow exponentially (exploding gradients).

  Mitigation strategies:
  \begin{itemize}
  \item gradient clipping
  \item orthogonal/unitary \(W_{hh}\)
  \item gating (LSTM/GRU)
  \end{itemize}
\end{frame}

%---------------------------------------------------------
\begin{frame}{BPTT Algorithm (Practical Implementation)}
  Forward pass (store \(a_t,h_t,o_t,y_t\) for \(t=1,\dots,T\)).

  Backward pass:
  \begin{enumerate}
  \item Initialize \(\delta^a_{T+1}=0\).
  \item For \(t=T,\dots,1\):
    \begin{itemize}
    \item compute \(\delta^o_t=\frac{\partial \ell_t}{\partial y_t}\odot\phi'(o_t)\)
    \item compute \(\delta^h_t = W_{hy}^T\delta^o_t + W_{hh}^T\delta^a_{t+1}\)
    \item compute \(\delta^a_t = \delta^h_t \odot \sigma'(a_t)\)
    \item accumulate gradients:
      \[
      \nabla W_{hy} {+}{=} \delta^o_t h_t^T,\quad
      \nabla b_y {+}{=} \delta^o_t
      \]
      \[
      \nabla W_{xh} {+}{=} \delta^a_t x_t^T,\quad
      \nabla W_{hh} {+}{=} \delta^a_t h_{t-1}^T,\quad
      \nabla b_h {+}{=} \delta^a_t
      \]
    \end{itemize}
  \end{enumerate}
\end{frame}

%---------------------------------------------------------
\begin{frame}{Truncated BPTT}
  For very long sequences, backpropagating through all \(T\) steps is costly and unstable.

  Truncated BPTT:
  \begin{itemize}
  \item Split sequence into blocks of length \(K\).
  \item Backpropagate only \(K\) steps at a time.
  \end{itemize}

  Effect:
  \begin{itemize}
  \item reduces cost from \(O(T)\) memory to \(O(K)\)
  \item limits long-range credit assignment
  \end{itemize}
\end{frame}


%------------------------------------------------
\section{Long Short-Term Memory}

\begin{frame}{Motivation for LSTM}

Simple RNNs struggle with long-range dependencies.

LSTMs introduce gates that control information flow.
\end{frame}

%------------------------------------------------
\begin{frame}{LSTM Equations}

Forget gate:

\[
f_t = \sigma(W_f x_t + U_f h_{t-1} + b_f)
\]

Input gate:

\[
i_t = \sigma(W_i x_t + U_i h_{t-1} + b_i)
\]

Candidate state:

\[
\tilde c_t = \tanh(W_c x_t + U_c h_{t-1})
\]
\end{frame}

%------------------------------------------------
\begin{frame}{Cell State Update}

\[
c_t = f_t \odot c_{t-1} + i_t \odot \tilde c_t
\]

Output gate:

\[
o_t = \sigma(W_o x_t + U_o h_{t-1})
\]

Hidden state:

\[
h_t = o_t \odot \tanh(c_t)
\]
\end{frame}

%------------------------------------------------
\section{Gated Recurrent Units}

\begin{frame}{GRU Architecture}

Simplified alternative to LSTM.

Update gate:

\[
z_t = \sigma(W_z x_t + U_z h_{t-1})
\]

Reset gate:

\[
r_t = \sigma(W_r x_t + U_r h_{t-1})
\]
\end{frame}

%------------------------------------------------
\begin{frame}{GRU Update}

\[
\tilde h_t =
\tanh(W_h x_t + U_h (r_t \odot h_{t-1}))
\]

\[
h_t =
(1-z_t)h_{t-1} + z_t \tilde h_t
\]
\end{frame}

%------------------------------------------------
\section{Computational Considerations}

\begin{frame}{Computational Complexity}

Cost per time step:

\[
O(n_h^2)
\]

where \(n_h\) is the hidden dimension.

Total cost:

\[
O(Tn_h^2)
\]
\end{frame}

%------------------------------------------------
\section{Applications in physical sciences}

\begin{frame}{RNN applications}

Examples:

\begin{itemize}
\item quantum state reconstruction
\item time-series modelling
\item turbulence prediction
\item particle detector signals
\item dynamical system forecasting
\end{itemize}
\end{frame}

%------------------------------------------------
\begin{frame}{Example: Dynamical Systems}

Given data

\[
x_{t+1} = F(x_t)
\]

RNN learns

\[
x_{t+1} \approx f_\theta(x_t,h_t)
\]

This provides a learned surrogate model.
\end{frame}

%------------------------------------------------
\begin{frame}{Summary}

\begin{itemize}
\item RNNs model sequential data using hidden states
\item training uses backpropagation through time
\item gradients can vanish or explode
\item LSTM and GRU mitigate these problems
\item widely used in physics and machine learning
\end{itemize}

\end{frame}

\end{document}



